% !TeX program = xelatex
\documentclass[a4paper,10pt]{article}

% ============================================================================
% ក្រដាសដាច់ដោយឡែក៖ សេណារីយ៉ូនៃការបង្រៀន (Scenario នៃការបង្រៀន – ៥០ នាទី)
% ============================================================================
\usepackage[a4paper, left=0.95cm, right=0.95cm, top=0.95cm, bottom=1.05cm, footskip=0.50cm]{geometry}
\usepackage{amsmath}
\usepackage{unicode-math}
\usepackage{fontspec}
\usepackage{array, longtable, colortbl}
\usepackage[table]{xcolor}
\usepackage{pifont}
\usepackage{fancyhdr}
\usepackage{tikz}

% LaTeX 2025 + longtable compatibility guard
\makeatletter
\let\@expl@@@mark@update@singlecol@structures@@\relax
\makeatother

% Cross-platform font path resolver (Windows local + Linux/Cloud repo fonts/)
\providecommand{\NGSFindFonts}{%
  \IfFileExists{C:/Windows/Fonts/cambria.ttc}{%
    \def\NGSSysPath{C:/Windows/Fonts/}%
    \def\NGSUserPath{C:/Users/KIT/AppData/Local/Microsoft/Windows/Fonts/}%
  }{%
    \IfFileExists{./fonts/KhmerOSsiemreap.ttf}{\def\NGSSysPath{./fonts/}\def\NGSUserPath{./fonts/}}{%
    \IfFileExists{../fonts/KhmerOSsiemreap.ttf}{\def\NGSSysPath{../fonts/}\def\NGSUserPath{../fonts/}}{%
    \IfFileExists{../../fonts/KhmerOSsiemreap.ttf}{\def\NGSSysPath{../../fonts/}\def\NGSUserPath{../../fonts/}}{%
    \IfFileExists{../../../fonts/KhmerOSsiemreap.ttf}{\def\NGSSysPath{../../../fonts/}\def\NGSUserPath{../../../fonts/}}{%
    \IfFileExists{../../../../fonts/KhmerOSsiemreap.ttf}{\def\NGSSysPath{../../../../fonts/}\def\NGSUserPath{../../../../fonts/}}{%
    \IfFileExists{../../../../../fonts/KhmerOSsiemreap.ttf}{\def\NGSSysPath{../../../../../fonts/}\def\NGSUserPath{../../../../../fonts/}}{%
    \IfFileExists{../../../../../../fonts/KhmerOSsiemreap.ttf}{\def\NGSSysPath{../../../../../../fonts/}\def\NGSUserPath{../../../../../../fonts/}}{%
      \def\NGSSysPath{../../../../../../../fonts/}\def\NGSUserPath{../../../../../../../fonts/}%
    }}}}}}}%
  }%
}
\NGSFindFonts

\setmathfont[
  Path={\NGSSysPath},
  FontIndex=1,
  Scale=1.02
]{cambria.ttc}

\defaultfontfeatures{Ligatures=TeX}

\setmainfont[
  Path={\NGSUserPath},
  Extension=.ttf,
  Script=Khmer,
  BoldFont={KhmerOSsiemreap},
  BoldFeatures={FakeBold=1.8},
  ItalicFont={KhmerOSsiemreap},
  ItalicFeatures={FakeSlant=0.25},
  BoldItalicFont={KhmerOSsiemreap},
  BoldItalicFeatures={FakeBold=1.8, FakeSlant=0.25}
]{KhmerOSsiemreap}

\newfontfamily\muol[
  Path={\NGSUserPath},
  Extension=.ttf,
  Script=Khmer
]{KhmerOSmuollight}

\newfontfamily\TimesEn[
  Path={\NGSSysPath},
  Extension=.ttf,
  Scale=1.00,
  BoldFont={timesbd},
  ItalicFont={timesi},
  BoldItalicFont={timesbi}
]{times}

\newfontfamily\NagaBold[
  Path={\NGSUserPath},
  Extension=.ttf
]{AKbalthom Naga Bold}

\newfontfamily\CalibriFont[
  Path={\NGSSysPath},
  Extension=.ttf
]{calibri}

\XeTeXlinebreaklocale "km"
\XeTeXlinebreakskip = 0pt plus 1pt minus 0.1pt

\definecolor{stepgray}{HTML}{D9D9D9}
\definecolor{headgray}{HTML}{F2F2F2}
\definecolor{darkred}{HTML}{C00000}
\definecolor{red5e}{HTML}{FF0000}

\pagestyle{fancy}
\fancyhf{}
\renewcommand{\headrulewidth}{0pt}
\fancyfoot[R]{\CalibriFont\fontsize{11pt}{13pt}\selectfont\thepage}

\setlength{\parindent}{0pt}
\setlength{\parskip}{0pt}
\setlength{\tabcolsep}{4pt}

\newcolumntype{L}[1]{>{\raggedright\arraybackslash}p{#1}}
\newcommand{\en}[1]{{\TimesEn #1}}
\newcommand{\enbf}[1]{{\TimesEn\textbf{#1}}}
\newcommand{\fiveE}[1]{{\NagaBold\fontsize{9pt}{12.5pt}\selectfont\textcolor{red5e}{#1}}}
\newcommand{\ditem}{\tikz[baseline=-0.5ex]\fill (0,0.40ex) -- (0.40ex,0) -- (0,-0.40ex) -- (-0.40ex,0) -- cycle;\hspace{0.25em}}
\newcommand{\chkbox}{\tikz[baseline=-0.1ex]\draw[black, line width=0.75pt](0,0) rectangle (0.23,0.23);}
\newcommand{\cnum}[1]{{\fontsize{9.5pt}{11pt}\selectfont\ding{\numexpr181+#1\relax}}\hspace{0.15em}}
\newcommand{\tri}[1]{\ensuremath{\triangle #1}}
\newcommand{\ang}[1]{\ensuremath{\angle #1}}
\newcommand{\cg}{\ensuremath{\cong}}
\newcommand{\ncg}{\ensuremath{\not\cong}}

\begin{document}
\fontsize{9.4pt}{14.3pt}\selectfont

% -------------------- ក្បាលឯកសារ Scenario --------------------
\begin{center}
  {\muol\fontsize{11pt}{16pt}\selectfont ព្រះរាជាណាចក្រកម្ពុជា}\\[1pt]
  {\muol\fontsize{10pt}{15pt}\selectfont ជាតិ សាសនា ព្រះមហាក្សត្រ}\\[4pt]
  {\muol\fontsize{10.5pt}{15pt}\selectfont សេណារីយ៉ូនៃការបង្រៀន \enbf{(Scenario} នៃការបង្រៀន\enbf{)}}
\end{center}

\vspace{2pt}
{\renewcommand{\arraystretch}{1.06}
\begin{tabular}{@{}l l l p{0.5cm} l l l@{}}
  - អនុវិទ្យាល័យ & : & ព្រះក្រសាំង (កំពង់ស្វាយ-កំពង់ធំ) & & - មុខវិជ្ជា & : & គណិតវិទ្យា ថ្នាក់ទី៨ \\
  - មេរៀនទី១២ & : & ២.១ លក្ខណៈមេដ្យាទ័រនៃអង្កត់ & & - រយៈពេលសរុប & : & \textcolor{darkred}{\textbf{១ ម៉ោងសិក្សា (៥០ នាទី)}} \\
  - គ្រូបង្រៀន & : & ហុក ប៊ុនថែន & & - កាលបរិច្ឆេទ & : & ថ្ងៃអង្គារ ទី១ ខែធ្នូ ឆ្នាំ២០២៦ \\
\end{tabular}
}

\vspace{3pt}

% -------------------- I. តារាងបែងចែកពេលវេលា និងសេណារីយ៉ូលម្អិត (៥០ នាទី) --------------------
\textbf{I.} {\muol តារាងបែងចែកពេលវេលា និងសេណារីយ៉ូប្រតិបត្តិការក្នុងថ្នាក់} \textbf{(៥០ នាទី)}

\vspace{-2pt}
\renewcommand{\arraystretch}{1.04}
\begin{longtable}{|c|L{2.75cm}|c|L{7.65cm}|L{4.45cm}|}
\hline
\rowcolor{headgray}
\multicolumn{1}{|c|}{\muol នាទី} &
\multicolumn{1}{c|}{{\muol ជំហាន \enbf{/ 5Es}}} &
\multicolumn{1}{c|}{\muol រយៈពេល} &
\multicolumn{1}{c|}{{\muol សេណារីយ៉ូប្រតិបត្តិ \enbf{(}ពាក្យពេចន៍គ្រូ \enbf{\&} សម្ភារៈ\enbf{)}}} &
\multicolumn{1}{c|}{{\muol ការគ្រប់គ្រង \enbf{\&} វិធីបម្រុង}} \\
\hline
\endfirsthead

\hline
\rowcolor{headgray}
\multicolumn{1}{|c|}{\muol នាទី} &
\multicolumn{1}{c|}{{\muol ជំហាន \enbf{/ 5Es}}} &
\multicolumn{1}{c|}{\muol រយៈពេល} &
\multicolumn{1}{c|}{{\muol សេណារីយ៉ូប្រតិបត្តិ \enbf{(}ពាក្យពេចន៍គ្រូ \enbf{\&} សម្ភារៈ\enbf{)}}} &
\multicolumn{1}{c|}{{\muol ការគ្រប់គ្រង \enbf{\&} វិធីបម្រុង}} \\
\hline
\endhead

\hline
\endfoot

\hline
\endlastfoot

00--02 & \textbf{ជំហានទី១}\newline រដ្ឋបាលថ្នាក់ & ២ នាទី &
\ditem ពិនិត្យអវត្តមាន និងអនាម័យថ្នាក់រៀន ។\newline
\ditem ចាត់ក្រុម ៣ នាក់ («ក្រុមគណៈកម្មការវាស់ស្ទង់ទីតាំងស្មើចម្ងាយ») ចែកកញ្ចប់ ១:១ (ទំព័រ ១) ដាក់ផ្កាប់ទុកលើតុ ។ &
ចាត់ប្រធានក្រុមស្រាប់ ហាមកាត់ក្រដាសឬបត់មុនបញ្ជាគ្រូ ។ \\ \hline

02--05 & \textbf{ជំហានទី២ (ក)}\newline កំណែកិច្ចការផ្ទះ & ៣ នាទី &
\ditem \textbf{សំណួរបញ្ឆោត (\#12 ទំ.១៦៥)}៖ «ហេតុអ្វីក្នុងលំហាត់ទី១២ ពី $\ang{CBA}=\ang{DAB}$ យើងទាញបាន $BC \parallel AD$?»\newline
\ditem ផ្ទៀងខ្លី \#10 ($\tri{BDC}\cg\tri{CEB}$ អ.ជ) \& \#12 (អ.ជ $\to$ មុំឆ្លាស់ក្នុងស្មើគ្នា) ។ &
កុំឱ្យសិស្សសរសេរវែងលើក្ដារខៀន ផ្ដោតលើមុំឆ្លាស់ក្នុងស្មើគ្នា ។ \\ \hline

05--08 & \textbf{ជំហានទី២ (ខ)}\newline រំលឹកមេរៀនទាក់ទង & ៣ នាទី &
\ditem សួររហ័ស៖ (១) $I$ ចំណុចកណ្ដាល $BC$ ($\Rightarrow IB=IC$) និង (២) បន្ទាត់ $\Delta\perp BC$ ត្រង់ $I$ ($\Rightarrow \ang{MIB}=\ang{MIC}=90^\circ$) ។\newline
\ditem រំលឹកលក្ខខណ្ឌប៉ុនគ្នា ជ.ម.ជ (SAS) និង ជ.ជ.ជ (SSS) ។ &
គូសសញ្ញា \textbf{=} លើ $IB, IC$ និងប្រអប់ $90^\circ$ លើក្ដារខៀន ។ \\ \hline

08--11 & \fiveE{E1 Engage}\newline បូមទឹកសូឡា & ៣ នាទី &
\ditem លើករឿង \textbf{«បញ្ហាដាក់ស្ថានីយបូមទឹកសូឡាស្មើចម្ងាយពីភូមិ $B$ និងភូមិ $C$»} ($MB=MC$)៖ ចំណុចកណ្ដាល $I$ ស្ថិតក្នុង \textbf{ប្រឡាយទឹកលិច} ។\newline
\ditem ស្រង់មតិពីទីតាំងដីទួល $M$ ក្រៅ $I$ រួចសរសេរចំណងជើង ២. និង ២.១ ។ &
គូស $BC$ ទ្រេត $12^\circ$ លើក្ដារខៀន ចៀសវាងគិតថា $\Delta$ ត្រូវតែដេក ។ \\ \hline

11--22 & \fiveE{E2 Explore}\newline បត់ \& Dual POE & ១១ នាទី &
\ditem \textbf{សកម្មភាព ១ (៦ ន.)}៖ លើទំព័រ ១ (១:១, ផ្ទៀងក្រិត $5.0\text{cm}$) កាត់បន្ទះ ១ ($BC=6.0\text{cm}$) បត់ $B\mapsto C$ រកផ្នត់ $\Delta$ កាត់ $I$, វាស់ $IB=IC=3.0\text{cm}, 90^\circ$, $M$ ($2.5\text{cm}$ លើ $I$, $MB=MC\approx 3.9\text{cm}$) \& $N$ ($3.5\text{cm}$ ក្រោម $I$, $NB=NC\approx 4.6\text{cm}$) ។\newline
\ditem \textbf{សកម្មភាព ២ Dual POE (៥ ន.)}៖ \textbf{POE 2A} ផ្ទាំង X ($60^\circ$ ត្រង់ $I$: $P_1C=3.0\neq P_1B\approx 5.2\text{cm}$) vs ផ្ទាំង Y ($\perp$ ត្រង់ $H$, $HB=2.0, HC=4.0$: $Q_1B\approx 3.6\neq Q_1C=5.0\text{cm}$) $\to$ \textbf{សោភ្លោះ}; \textbf{POE 2B} ធ្នូដែកឈាន $r=4.5\text{cm}$ កាត់ត្រង់ $P, Q \Rightarrow (PQ)\equiv\Delta$ ។ &
បោះពុម្ពទំព័រ ១ ខ្នាត \textbf{Actual Size 100\%} ។ គ្រូដើរត្រួតសិស្សបត់ $B$ ត្រួតលើ $C$ ពិតប្រាកដ និងប្រើដែកឈាន $r=4.5\text{cm}$ ។ \\ \hline

22--28 & \fiveE{E3 Explain}\newline ផ្ទាល់ \& \cnum{1}--\cnum{5} & ៦ នាទី &
\ditem បង្គុំនិយមន័យ (ទំ.១៥៦) និងសម្រាយផ្ទាល់ \textbf{ជ.ម.ជ} លើរូបបញ្ឈរ $BC$ ($C$ លើ $B$ ក្រោម, $M$ ឆ្វេង \textbf{/}, $N$ ស្ដាំ $\times$)៖ សង្កត់ $IM$ រួម, \textbf{$=$ vs $\cg$}, និងវិបាក $MB=MC$ ។\newline
\ditem បិទផ្ទាំង \textbf{ធនាគារហេតុផល \cnum{1}--\cnum{5}} (ទំព័រ ៤) + កែអក្ខរាវិរុទ្ធទំ.១៥៦ (បំពេញ \textbf{«។»} ចុងបន្ទាត់ ១, សម្រួល \textbf{«ស្ថិតនៅ»}) + ញែក \textbf{ពុះមុំ vs មេដ្យាទ័រ} ។ &
ប្រើ ២ ចំណុចសម្រាយ៖ សោភ្លោះ ($IB=IC$ + $90^\circ$) និងភាពខុសគ្នាពីពុះមុំ (កែងជ្រុង) ។ \\ \hline

28--38 & \fiveE{E4 Elaborate}\newline ច្រាស \& ២ ចំណុច & ១០ នាទី &
\ditem ចែកកាត \textbf{«កាតសោភ្លោះ និងធ្នូប្រសព្វពីរចំណុច»} (ទំព័រ ២)៖\newline
+ ២.៥ន. គូសបន្ទាត់ជំនួយ $MI$ ដាច់ៗ រួច \textbf{[tracing]} លើ \textbf{/}, $\times$, \textbf{=}\newline
+ ៥.៥ន. សរសេរសម្រាយ៖ (ក) ទ្រឹស្ដីច្រាសតាម \textbf{ជ.ជ.ជ} ($\tri{MIB}\cg\tri{MIC}\Rightarrow 90^\circ$) \& (ខ) ២ ចំណុច $MB=MC, NB=NC \Rightarrow (MN)\equiv\Delta$ ដោយមិនបាច់វាស់ $I$\newline
+ ២ន. \textbf{EIA 30s}៖ អន្ទាក់ ១ ចំណុច ($PB=PC$ មានបន្ទាត់ទ្រេតច្រើន) vs ២ ចំណុច ។ &
ដាស់តឿនមុនសរសេរ៖ ហាមសន្មត $\ang{MIB}=90^\circ$ មុនស្រាយ (ក) ត្រូវទាញពីមុំរាប $180^\circ$! \\ \hline

38--48 & \fiveE{E5 Evaluate}\newline ប្រតិបត្តិ \& កាត & ១០ នាទី &
\ditem \textbf{ប្រតិបត្តិ «គណៈកម្មការផ្ទៀងសោភ្លោះ ៣ តួនាទី» (៦.៥ ន.)} (ទំព័រ ៣, $BC=8\text{cm}, MB=3x-1, MC=x+7, NB=NC=6\text{cm}$)៖\newline
+ \textbf{តួនាទី ១ (Pause-Ask)}៖ ផ្ទៀងសោភ្លោះ $I$ និងហេតុអ្វី $MB=MC$ (\cnum{4})\newline
+ \textbf{តួនាទី ២}៖ គណនា $IB=IC=4, x=4 \Rightarrow MB=MC=11\text{cm}$ (\cnum{1}--\cnum{4})\newline
+ \textbf{តួនាទី ៣ (Pause-Review)}៖ ស្រាយ \cnum{5} ($N\in\Delta\Rightarrow M,I,N$ ត្រង់ជួរ) + \textbf{បិទ E1} ។\newline
\ditem \textbf{កាតត្រួតពិនិត្យ A--D (៣.៥ ន.)} (ទំព័រ ៥)៖ A (\ding{51} $MB=5$), B (\ding{55} $65^\circ$), C (\ding{55} $2\neq 5$), D (\ding{55} ១ ចំណុច) + ម្រាម ១--៣ ។ &
បើពេលជិតអស់ បង្រួមកាត A--D មក ៩០ វិនាទី ដោយសួរមូលហេតុខ្វះសោមួយណា (B/C/D) ។ \\ \hline

48--50 & \textbf{ជំហានទី៥}\newline កិច្ចការផ្ទះ \& បណ្ដាំ & ២ នាទី &
\ditem ចែកប័ណ្ណកិច្ចការផ្ទះ (ទំព័រ ៦, ទំ.១៦៤)៖ \#1 ($BC=2x+30, EF=5x-90$), \#3 ($MT=MB, MA=MO$ បំបែកត្រួត), \#5 (ក. $AB\parallel CD, TR=TS$; ខ. $\Delta\perp RS$ កាត់ $T$) ។\newline
\ditem ណែនាំអានមុន §២.២ (ទំ.១៥៧--១៥៨) និងសុវត្ថិភាពចរាចរណ៍ ។ &
ផ្ញើរូបបង្ហាញលំហាត់ \#1, \#3, \#5 ចូល \en{Telegram} (បញ្ជាក់ \#2 ប្ដូរទៅ ២.២) ។ \\ \hline

\rowcolor{stepgray}
\multicolumn{2}{|r|}{\textbf{សរុបរយៈពេលបង្រៀនពេញ}} & \textbf{៥០ នាទី} &
\multicolumn{2}{l|}{ត្រូវគ្នា ១០០\% នឹង \en{LS\_Plan\_PB.tex} និង \en{Printable\_Materials\_PB.tex} (ទំព័រ ១--៦)} \\ \hline
\end{longtable}

\vspace{2pt}

% -------------------- II. បញ្ជីត្រៀមសម្ភារៈ & III. ការគ្រប់គ្រងកំហុស --------------------
\begin{minipage}[t]{0.46\textwidth}
  \textbf{II.} {\muol បញ្ជីត្រៀមសម្ភារៈមុនចូលថ្នាក់} \textbf{(ស-ខ្មៅ)}\\[3pt]
  \begin{tabular}{|c|L{7.2cm}|}
    \hline
    \rowcolor{headgray}
    \textbf{\ding{51}} & \multicolumn{1}{c|}{\textbf{សម្ភារៈត្រូវរៀបចំទុកជាមុន}} \\ \hline
    \chkbox & កញ្ចប់ ១:១ ទំព័រ ១ (\textbf{Actual Size 100\%}, ក្រិត $5.0\text{cm}$)៖ បន្ទះបត់ $BC=6.0\text{cm}$, ផ្ទាំង X/Y (POE 2A) \& ធ្នូ $r=4.5\text{cm}$ (POE 2B) \\ \hline
    \chkbox & កន្ត្រៃ បន្ទាត់ក្រិត ១:១ រ៉ាព័រទ័រ និងដែកឈាន \\ \hline
    \chkbox & កាតសោភ្លោះ E4 (ទំព័រ ២), សន្លឹក ៣ តួនាទី E5 (ទំព័រ ៣) \& ផ្ទាំង \cnum{1}--\cnum{5} (ទំព័រ ៤) \\ \hline
    \chkbox & កាត Exit \textbf{A--D} (ទំព័រ ៥) \& ប័ណ្ណកិច្ចការផ្ទះ \#1, \#3, \#5 (ទំព័រ ៦) \\ \hline
  \end{tabular}
\end{minipage}%
\hfill
\begin{minipage}[t]{0.51\textwidth}
  \textbf{III.} {\muol កំហុសភ័ន្តច្រឡំ} \textbf{(តាមការស្រាវជ្រាវ) \&} {\muol ដំណោះស្រាយ}\\[3pt]
  \begin{tabular}{|L{8.8cm}|}
    \hline
    \rowcolor{headgray}
    \multicolumn{1}{|c|}{\textbf{ឧបសគ្គសិស្ស (Empirical Misconceptions) \& វិធានការទប់ស្កាត់}} \\ \hline
    \textbf{១. អគតិសោម្ខាង (Partial-Lock) \& រូបបញ្ឈរ៖} បន្ទាត់ទ្រេត $60^\circ$ (X) vs កែងខុសកណ្ដាល (Y) និងបត់បន្ទះទ្រេត $12^\circ$ \en{(Hodancı \& Hacısalihoğlu Karadeniz, 2025; Catman-Aksoy \& Isiksal-Bostan, 2024)} ។ \\ \hline
    \textbf{២. ច្រឡំពុះមុំ vs មេដ្យាទ័រ \& Locus ភូមិពីរ៖} បន្ទាត់ ២ ចម្ងាយកែងជ្រុង vs ចុងអង្កត់ និងបរិបទស្ថានីយសូឡា \en{(Hodancı \& Hacısalihoğlu Karadeniz, 2025; Bento et al., 2018)} ។ \\ \hline
    \textbf{៣. បោះបង់ធ្នូដែកឈាន \& បន្ទាត់ជំនួយ $MI$៖} ប្រៀបធៀបកាំ $r=4.5\text{cm}$ ($>3\text{cm}$) ជំរុញគូស $MI$ ដាច់ៗ និង \textbf{[tracing]} \en{(Šipuš et al., 2022; Chen et al., 2017; Muzaini et al., 2023; Yeo \& Tzeng, 2019)} ។ \\ \hline
    \textbf{៤. រំលងជ្រុងរួម $IM=IM$, អន្ទាក់ ១ ចំណុច \& រូបត្រួតគ្នា ($53\%$)៖} ៣ តួនាទី \en{Pause-Ask}, \textbf{កាត D}, និងបំបែក $\tri{MTO}, \tri{MBA}$ (\#3) \en{(Winer \& Battista, 2022; Biber, 2020; Miyazaki et al., 2019)} ។ \\ \hline
  \end{tabular}
\end{minipage}

\end{document}
